\section{AegisKV Design}
\label{sec:design}

AegisKV separates complete recovery coverage from fast recovery access.
Shared DRAM and storage retain complete checkpoints; Peer HBM holds
selected hot blocks on a preassigned takeover DP. Placement accounts for
both workflow completion time (WCT) and normal cache reuse.
We target fail-stop compute-worker failures with a compatible surviving
DP. The directory, ordered token log and advertised backing must survive
the protected failure domain. The application exposes known dependencies
and accepted inputs; external tool side effects are outside the KV
recovery contract.

\subsection{Observations}
\label{sec:design-insights}

\paragraph{1. Recovery urgency and reuse value differ.}
A faster restore shortens WCT only if it advances a required path
(Figure~\ref{fig:motivation-recovery-exposure}). Yet a high-slack Agent
can reuse its history after a tool call; repeatedly reconstructing that
history consumes shared serving capacity. Retaining a complete copy and
paying for a faster copy are therefore separate decisions.

\paragraph{2. Sparse access creates selective-residency opportunities.}
Figure~\ref{fig:sparse-releases} surveys selected long-context releases
with block, token or compressed-entry selection. It motivates supporting
selective KV access, without implying a monotonic increase in sparsity.
In the native traces of Figure~\ref{fig:design-insights}, the hottest
approximately 20\% of historical blocks account for 87.64\%, 74.76\%
and 56.54\% of accesses across the three models. This retrospective
concentration motivates hot replicas; it does not establish future-query
prediction accuracy. Complete backing remains necessary for misses.

\begin{figure}[t]
 \centering
 \includegraphics[width=\columnwidth]{figures/sparse-models.pdf}
 \caption{Selected sparse/hybrid model releases and reported context
 lengths. IDs: 1 MiniCPM4~\cite{minicpm4model};
 2 MiniCPM4.1~\cite{openbmb2025minicpm41};
 3 DeepSeek-V3.2-Exp~\cite{deepseek2025v32exp};
 4 DeepSeek-V3.2~\cite{deepseek2025v32model};
 5 GLM-5~\cite{zai2026glm5model};
 6 MiniCPM-SALA~\cite{openbmb2026sala};
 7 DeepSeek-V4 preview~\cite{deepseek2026v4model};
 8 MiniMax-M3~\cite{minimax2026m3};
 9 GLM-5.2~\cite{zai2026glm52};
 10 LongCat-Flash-Lite-Sparse~\cite{longcat2026sparse}.
 Dates are announcements, or first official repository dates for
 IDs 5, 9 and 10. Markers identify the sparse branch's selection unit,
 including SALA's block branch and V4's compressed sparse branch;
 they do not imply uniform attention across layers.}
 \label{fig:sparse-releases}
 \Description{Ten selected sparse or hybrid attention releases between
 June 2025 and July 2026, with reported contexts from 32K to 1M.
 Circles, squares and diamonds distinguish block, token and compressed
 entry selection. This is a literature survey, not a sparsity-rate curve.}
\end{figure}

\begin{figure}[t]
 \centering
 \includegraphics[width=\columnwidth]{figures/kv-access-concentration.pdf}
 \caption{Native KV access concentration for NOSA-8B~\cite{huang2026nosav2},
 NSA-8B (MiniCPM-4-8B-NSA)~\cite{fasa2026minicpm4nsa}, and
 MiniCPM4.1-8B~\cite{openbmb2025minicpm41}. We use the final 512 queries
 of matched 32,704-token prefixes from three PG19 texts. Each query
 counts a historical 64-token block once across heads. Rank the 503
 blocks within each layer/window; average access shares over all
 32 layers, then texts. Bands show text min--max; the diagonal denotes
 uniform access. Compressed/selector state is excluded. This is
 retrospective concentration, not prediction accuracy or recovery speedup.}
 \label{fig:design-insights}
 \Description{Cumulative historical block access curves for NOSA-8B,
 NSA-8B and MiniCPM4.1-8B, averaged over every layer and three texts.
 About 20 percent of blocks account for 87.64, 74.76 and 56.54 percent
 of accesses, respectively; all three curves exceed the uniform baseline.}
\end{figure}

\paragraph{3. Only exposed transfers justify Peer residency.}
DRAM avoids historical reconstruction; Peer additionally removes lower-tier
read/H2D waiting. Later-layer transfers may already overlap computation
that must still execute. Useful Peer depth therefore depends on a
request's committed history, uncommitted tail and recovery queues.

\paragraph{4. Reuse value persists across requests but changes over time.}
Tool waits preserve an Agent's future interest, and another Agent can
reuse the same prefix. A one-time score misses these changes; renewing
credit on every checkpoint can instead preserve stale copies indefinitely.
We distinguish actual reuse from value updates, using renewable credit
and byte-proportional pressure aging inspired by Landlord and
GreedyDual-Size~\cite{young2002landlord,cao1997greedydual}.
A complete-plan check handles dependencies and movement; this adaptation
inherits no competitive guarantee.

\subsection{System Overview}
\label{sec:system-overview}

Figure~\ref{fig:system-overview} has three layers. The Application supplies
requests and workflow progress. In the Controller, \textbf{(1)} Tracker
estimates remaining time and slack, \textbf{(2)} Router binds execution
and takeover DPs, and \textbf{(3)} KV Master chooses complete backing and
Peer hot replicas. In each Engine, \textbf{(4)} Adapter controls the Peer
budget and \textbf{(5)} Connector executes transfers. Local execution,
Peer replicas and transient workspace share HBM with separate accounting.

\begin{figure}[t]
 \centering
 \includegraphics[width=\columnwidth]{figures/system-overview.pdf}
 \caption{AegisKV architecture. Callouts (1)--(5) match the text;
 DP1 mirrors DP0's Adapter/Connector. DP0 executes; DP1 hosts Peer
 replicas and takes over. Blue links carry KV, dark dashes control,
 red dashes retry; orange blocks are unsealed tails. Complete shared
 backing lies outside both DP failure boundaries.}
 \label{fig:system-overview}
 \Description{Compact Application, Controller and Engine layers with
 numbered components, symmetric DP Local/Peer HBM and shared backing.}
\end{figure}

The storage object is an immutable, version-compatible token block in
one layer and TP shard. A complete checkpoint manifest resolves all layers,
shards, required selector state and the matching log prefix through
frontier $M$, entirely in DRAM or entirely in storage. Objects can be
shared across manifests; each physical replica consumes capacity once.
Sparse selection units, engine blocks and kernel tiles need not coincide;
the 64-token block in our measurements is not a universal requirement.

Prefill offers sealed state; decode triggers review every configured,
block-aligned $Z$ tokens. Arrival, tool handoff, demand/slack changes and
pressure also trigger review. A final snapshot seals a short round's
partial block using valid length and copy-on-write; accepted transfers
hold source leases until completion. Only a complete commit advances
$M$, so the logged tail $N$ can exceed $Z$ when copying lags. Peer coverage
requires complete committed backing.

\subsection{Workflow Tracker: Remaining Time and Slack}
\label{sec:workflow-protection}

Tracker maintains the residual workflow DAG. Each node is an Agent's
sequential segment between exposed fork/join boundaries or completion,
possibly spanning several LLM calls and tool waits. Child execution is
represented on child paths, not counted again in the parent; unknown
future forks are excluded.

A GBDT predicts segment seconds remaining from causal features: role/model,
current input length, elapsed time, completed-call count and previous
calls' lengths, timings and context growth. Missing history is masked;
training, calibration and evaluation split whole workflows. Labels end
at the segment boundary, and future call outcomes are never features.
Forward/backward critical-path passes~\cite{kelley1959cpm} derive slack
(Algorithm~\ref{alg:workflow-tracker}). Requests in one segment share
its margin; the Master evaluates their joint DAG effect.

\par\medskip\noindent\begin{minipage}{\columnwidth}
\raggedright
 \refstepcounter{designalgorithm}\label{alg:workflow-tracker}
 \textbf{Algorithm \thedesignalgorithm: Remaining time and slack}
 \par\smallskip
 \begin{tabularx}{\columnwidth}{@{}r@{\quad}X@{}}
 \toprule
 1 & $(G,X)\gets\textsc{SnapshotProgress}()$ \\
 2 & $\nu\gets\textsc{NextVersion}()$ \\
 3 & \textbf{if not} \textsc{Supported}$(G,X)$ \textbf{then} \\
 4 & \quad \textsc{PublishInvalid}$(\nu)$; \textbf{return} \\
 5 & $d\gets\max(0,\textsc{GBDT}(X));\quad g\gets\textsc{Guard}()$ \\
 6 & $EF\gets\textsc{Forward}(G,d)$ \\
 7 & $LF\gets\textsc{Backward}(G,d,\max_v EF_v)$ \\
 8 & \textbf{for} $v\in V(G)$ \textbf{do} \\
 9 & \quad $s_v\gets\max(0,LF_v-EF_v-g_v)$ \\
 10 & \textsc{Publish}$(\nu,G,d,g,s)$ \\
 \bottomrule
 \end{tabularx}
\end{minipage}\par\medskip

Here $G$ is the residual DAG; $X$ contains features; $d$, $g$ and $s$ are
per-segment duration, guard and advisory slack. The guard is an upper
empirical quantile of positive slack overestimation on calibration
workflows, not a per-request guarantee. \textsc{Supported} checks input
support, dependencies, calibration and freshness. Invalid estimates grant
no positive slack credit. Progress events and periodic refresh publish
versioned updates, invalidating old placement costs. The DAG passes cost
$O(|V|+|E|)$; prediction quality is evaluated separately in
Section~\ref{sec:eval-prediction}.

\subsection{Recovery-aware Router}
\label{sec:recovery-router}

Normal dispatch balances load and compatible prefix locality.
Figure~\ref{fig:agent-prefix-reuse} reports 97.90\% potential same-agent
token coverage and 0.37\% additional cross-agent coverage. Agent affinity
is a preference: shared prefixes can still help another Agent's first call.
Tool handoff releases the request's execution references and logical Peer
pin while retaining reusable backing.

\begin{figure}[t]
 \centering
 \includegraphics[width=\columnwidth]{figures/agent-prefix-reuse.pdf}
 \caption{Potential prefix reuse over 98,827 requests, 2,090 agents and
 393 CC/Weka workflows. Coverage uses earlier completed requests with
 the same recorded model and workflow. Same/other-agent coverage
 overlaps; only additional cross-agent coverage extends same-agent
 reuse. These are token-prefix opportunities, not runtime cache hits.}
 \label{fig:agent-prefix-reuse}
 \Description{Same-agent prefixes dominate aggregate reuse; cross-agent
 history can also help first subagent calls.}
\end{figure}

At arrival, Router binds a versioned, layout-compatible takeover DP outside
the source failure domain; its corresponding TP shards host Peer copies.
On confirmed failure, Router fences the old epoch and activates that target
subject to workspace and admission constraints. Rebinding preserves a valid
backing path and starts with zero assumed Peer coverage on the new target.
Before output acknowledgement, a surviving log records request identity,
token position/value and epoch; prompt and accepted tool inputs remain
recoverable. Replay suppresses duplicate outputs and stale callbacks.

\subsection{KV Master: Three-Tier Renewable Credit}
\label{sec:kv-master}

\subsubsection{Valuing complete backing and Peer acceleration}
\label{sec:placement-value}

At a checkpoint event, Master proposes a complete DRAM extension, then
storage if DRAM cannot admit it. It also considers complete promotion or
demotion, last-copy removal with an older/logged fallback, and Peer hot-set
replacement. Unchanged objects are shared. If neither backing tier admits
an extension, $M$ stays unchanged; large slack alone does not suppress
normal-reuse retention.

\paragraph{Tier-conditional value.}
Storage buys retention above reconstruction or an older checkpoint.
DRAM buys faster access above complete storage; without storage it also
provides retention. Peer buys only exposed read/H2D wait above complete
backing, never the recomputation saving already provided there. Demotion
must secure finite lower-tier capacity and refresh dependent Peer values.
Temporary duplicates do not receive duplicate retention value.

\paragraph{Prefix credit.}
Fix an Agent's known history and tail. Starting from the tier's actual
fallback, add legal groups in prefix order: complete token extents across
all layers, TP shards and auxiliary state for backing; layer/hot-block
groups for Peer. Each step can retain its previous restore path, making
local gains nonnegative. Each group receives the expected incremental reduction
in normal-use service time plus recovery delay beyond guarded slack.
Normal reuse is not clipped by slack, and an existing local hit earns no
extra backing benefit. Forecasts use visible reuse history and the same
finite horizon and scenarios as the plan model below.

Split each increment across its new physical objects by bytes. Distinct
Agents contribute separate shares to a shared prefix; repeated references
by one Agent do not duplicate its share. For example, successive restore
times 12, 7 and 5\,s with slack 6\,s yield failure-value increments 5 and
1\,s before probability weighting, rather than assigning all 6\,s to the
tail. Increments telescope under the Agent's local model; their sum across
Agents is a screening proxy, not exact common-failure WCT.

\begin{figure}[t]
 \centering
 \includegraphics[width=\columnwidth]{figures/retention-value.pdf}
 \caption{Slack alone does not rank retention value. In this isolated-fault
 example, A uses two capacity units to reduce recovery from 4 to 2\,s;
 B uses one to reduce 18 to 8\,s. A at slack 0 saves 1\,s/unit, whereas
 B at slack 6 saves 10\,s/unit. Normal-use value is zero. These constructed
 profiles motivate cost/size comparison, not a policy-speedup claim.}
 \label{fig:retention-value}
 \Description{A higher-slack object can save more workflow time per byte.}
\end{figure}

\paragraph{Complete-plan check.}
For current placement $P$ and candidate $P'$, including victims and moves,
Master evaluates
\begin{equation}
 \Delta J=\sum_e p_e\sum_w[T_w(P',e)-T_w(P,e)].
 \label{eq:value}
\end{equation}
$T_w$ is predicted remaining workflow completion time. A fixed finite
scenario table covers normal continuation and source-domain faults at
specified offsets, with probabilities summing to one. Common faults
include every interrupted request. The resource model uses the same queue
snapshot, service profiles, finite workspace and Connector service order
for both plans; completed copies determine $M$ and fault-time $N$.
Its delays enter the residual DAG relative to the same ordinary-execution
reference. Joins take the maximum; copy interference and temporary
occupancy are counted in that schedule, not charged twice.

The guard is added once per interrupted segment in both plans before DAG
evaluation. Invalid workflow timing falls back to summed normal/recovery
request completion delays without slack credit; invalid resource profiles
pause optional moves and use bounded recency order for safe mandatory
relief. Thus accelerating one of two equally delayed join
branches need not improve WCT, even if its local credit is positive.

\subsubsection{Choosing per-request layers and hot blocks}
\label{sec:sparse-placement}

For request $r$, $L_r$ is the Peer prefix depth and $K_{r,l}$ the number
of token blocks in layer $l\le L_r$, including required TP objects.
Mandatory native dependencies come first; remaining blocks are ranked
by selection frequency over completed queries, counting each physical
block once per query across heads. Ties favor recency, then block ID.
A budget below the mandatory set is infeasible. Complete backing supplies
misses without changing native attention selections
(Figure~\ref{fig:sparse-recovery}).

\begin{figure}[t]
 \centering
 \includegraphics[width=\columnwidth]{figures/sparse-recovery.pdf}
 \caption{Complete backing, per-layer Peer budgets and demand misses.
 (a) The DRAM manifest covers B0--B3 in every layer; the tail is uncommitted.
 (b) On the bound target, $L=2$, with $K_1=2$ and $K_2=3$; layer 3 has
 no Peer copy. (c) Layer 1 needs B0/B2: B0 hits Peer and B2 is fetched
 from backing before computation. B3 was unused. The complete manifest
 also covers TP shards and selector state. If demoted as a whole,
 shared storage supplies the miss path before H2D.}
 \label{fig:sparse-recovery}
 \Description{Three complete historical layers, an uncommitted tail,
 different hot-block counts and an explicitly routed layer-1 miss.}
\end{figure}

Master enumerates a finite set of per-layer block-count or fraction
templates and prefix depths from zero to model depth, plus the current
placement. This avoids a Cartesian product of all layer budgets. Bounded
bundles can jointly advance tied join branches; all evaluations share one
work budget. Each scenario uses current $M$ and its own predicted $N$,
including commit lag. Unsupported sparse profiles use full demand.
Only computation that must still execute hides transfer.

Figure~\ref{fig:design-pipeline} isolates Peer's benefit: (c) and (d) share
complete DRAM backing and the same replay tail; Peer removes 24\,ms of
exposed waiting in this analytical instance. Its seven-layer depth is
profile-specific. Neither $L$ nor $K$ must increase with every token:
a fixed budget can change block identities, and $L=0$ is valid when no
exposed wait is removed. Proposed sets become resident only after commit.

\begin{figure}[t]
 \centering
 \includegraphics[width=\columnwidth]{figures/recovery-pipeline.pdf}
 \caption{Analytical layered recovery with and without Peer.
 MiniCPM4-8B has 32 layers, 32K history and a
 64-token tail. Replay is calibrated at 2.870\,ms/layer; assumed effective
 H2D bandwidth is 6.25\,GB/s (5.369\,ms per full layer). Trace-derived
 mean preload and preload-plus-miss fractions are 0.6584/0.7748.
 (a) waits for bulk restoration; (b) overlaps full-layer restoration;
 (c) preloads hot blocks and fetches revealed misses; (d) adds 147.5\,MiB
 of Peer hot sets in the first seven layers. Numbers are tail completion
 and exposed wait. Cold copies follow replay; the ideal preemptible
 stream and homogeneous layer costs omit bounded-workspace and queue
 effects. This is not end-to-end fault-recovery measurement.}
 \label{fig:design-pipeline}
 \Description{The same required replay computation with progressively
 reduced exposed KV waits: 172, 83, 44 and 20 milliseconds.}
\end{figure}

\subsubsection{Online renewal, admission, and eviction}
\label{sec:placement-plan}

A physical replica $x$ has size $b_x$, each Agent's current grant $g_{a,x}$
and paid rent $d_{a,x}$. Each pool has its own cumulative rent $\rho_j$;
$L_r$ denotes only layer depth. Remaining credit and the eviction key are
\begin{equation}
 C_x=\sum_a[g_{a,x}-d_{a,x}]^+,
 \qquad H_x=\rho_j+C_x/b_x.
 \label{eq:rental-credit}
\end{equation}
Clocks and new paid amounts start at zero. Advancing rent by $\delta$
charges every committed replica $\min(C_x,b_x\delta)$, split proportionally
among its positive Agent shares and added to their paid amounts. Pinned
and leased copies also age; eligibility for removal is separate. Settlement
can be lazy. Rent represents capacity competition, not wall time.

An actual prefix/recovery use settles rent and renews only that Agent's
share by refreshing its grant and resetting its paid amount. Request/use
IDs prevent repeated renewal within one use. Value-only changes refresh
grants but preserve paid rent; checkpoint growth does not renew unchanged
prefixes. Expired demand removes its grant, not the copy. Figure~\ref{fig:renewable-credit}
illustrates the distinction.

\begin{figure}[t]
 \centering
 \includegraphics[width=\columnwidth]{figures/renewable-credit.pdf}
 \caption{Constructed credit accounting for one shared physical object
 of unit size. Initial Agent grants are 3 and 2. A rent charge of 2 leaves
 shares 1.8 and 1.2. Updating A's grant to 4 preserves its paid amount
 1.2, leaving 2.8; an actual new use by A then renews its share to 4.
 B's share remains 1.2 in both updates. Credit is a screening quantity,
 not measured WCT or duplicated physical capacity.}
 \label{fig:renewable-credit}
 \Description{Two Agents contribute to one physical prefix object.
 A four-row ledger separates pressure aging, value refresh and renewal.}
\end{figure}

\paragraph{Legal competition.}
Victims are legal moves: complete-checkpoint demotion or fallback,
or valid Peer hot-set reduction. Bytes remain occupied while another
manifest references them; bounded owner bundles can release shared objects.
Ordinary competition respects running requests' logical Peer pins.
Request end/tool handoff unpins; self replacement and required capacity
relief may override logical pins, never physical read/copy leases.

For each move, \textsc{RentPlans} refreshes up to $q\ge2$ low-key victims
and $q$ rotating-scan alternatives, exploring victim chains within a common
work budget. Releasing objects $Y$ raises the trial clock to at least
$\max_{x\in Y}H_x$. A group of new bytes must retain positive credit:
$G_{\rm in}>b_{\rm in}(\rho'_j-\rho_j)$. The trial charge is split by
incoming grants; existing replicas neither consume new capacity nor gain
new credit. Each pool applies its own gate. Plans include transient space
and recursively secure storage for DRAM demotion; Peer release uses its
existing backing without redundant D2H. These bounded alternatives do
not guarantee the best feasible replacement.

\par\medskip\noindent\begin{minipage}{\columnwidth}
\raggedright
 \refstepcounter{designalgorithm}\label{alg:placement}
 \textbf{Algorithm \thedesignalgorithm: Three-tier rental event loop}
 \par\smallskip
 \begin{tabularx}{\columnwidth}{@{}r@{\quad}X@{}}
 \toprule
 1 & \textsc{Reconcile}$(e)$; \textsc{SettleAffected}$(e)$ \\
 2 & \textbf{if} \textsc{NewUse}$(e)$ \textbf{then} \textsc{RenewAgent}$(e)$ \\
 3 & \textsc{MarkChangedValues}$(e)$; $S\gets\textsc{Snapshot}()$ \\
 4 & $h\gets\textsc{ShrinkRequired}(S)$; $(P^*,d^*)\gets(\varnothing,\infty)$ \\
 5 & \textbf{for} $m\in\textsc{Candidates}(S,h)$ \textbf{while} \textsc{WorkLeft}() \\
 6 & \quad \textbf{for} $P\in\textsc{RentPlans}(S,m,h)$ \\
 7 & \qquad $d\gets\Delta J(P;S)$ \\
 8 & \qquad \textbf{if} $d<d^*$ \textbf{then} $(P^*,d^*)\gets(P,d)$ \\
 9 & \textbf{if} $P^*=\varnothing$ \textbf{then} \textsc{Defer}$(e)$; \textbf{return} \\
 10 & \textbf{if} $\neg h\land d^*\ge-\epsilon$ \textbf{then return} \\
 11 & \textbf{if not} \textsc{ReserveIfCurrent}$(S,P^*)$ \textbf{then} \\
 12 & \quad \textsc{Requeue}$(e)$; \textbf{return} \\
 13 & \textsc{Submit}$(P^*)$ \quad\textit{completion/failure emits a new event} \\
 \bottomrule
 \end{tabularx}
\end{minipage}\par\medskip

Algorithm~\ref{alg:placement} handles use, checkpoint, value, request-end,
budget, failure and transfer events. \textsc{Reconcile} applies verified
commits, unpins and failures; value refresh preserves paid rent. Optional
plans must pass rent screening and $\Delta J<-\epsilon$. Required relief
bypasses positive-credit admission and chooses the least costly examined
legal plan, still charging released objects. No safe plan pauses admission
and defers work; later events retry the newest pending frontier.
Periodic budgeted refresh revisits stale keys. Tool waits retain Agent
interest; terminal progress removes failure demand, and expired inactive
shares are pruned.

\textsc{ReserveIfCurrent} rechecks versions, leases, bytes and pool clocks.
Rejected previews change no committed rent or residency. A prepared
handoff advances rent only after its fallback is published; new replicas
install their residual trial credit on commit. Failure before handoff
discards trial rent, whereas failure after handoff cannot undo released
victims or reduce the clock. Connector completion/failure feeds the next
event, closing the loop without treating a pending copy as protection.

\subsection{Adapter: Guarded AIMD for Peer Capacity}
\label{sec:adapter}

Adapter controls total Peer authorization using pressure-driven
multiplicative decrease and stable-window additive
increase~\cite{chiu1989aimd}. Its physical ceiling is
\begin{equation}
 C=\max\{0,\min(B_{\mathrm{cap}},
 B_{\mathrm{pool}}-B_{\mathrm{serve}}-B_{\mathrm{reserve}})\}.
 \label{eq:peer-ceiling}
\end{equation}
The pool excludes weights/fixed allocations; serving high watermark
excludes Peer, and reserve covers growth/workspace. Existing Peer bytes
are not subtracted twice. $B_{\rm cap}$ can be a configured pool fraction.

\par\medskip\noindent\begin{minipage}{\columnwidth}
\raggedright
 \refstepcounter{designalgorithm}\label{alg:adapter}
 \textbf{Algorithm \thedesignalgorithm: Windowed Peer quota}
 \par\smallskip
 \begin{tabularx}{\columnwidth}{@{}r@{\quad}X@{}}
 \toprule
 1 & $C\gets\textsc{SafeCeiling}();\quad B'\gets B$ \\
 2 & \textbf{if} \textsc{Pressure}() \textbf{then} $(B',z)\gets(\gamma B,0)$ \\
 3 & \textbf{else if} \textsc{Stable}() \textbf{then} \\
 4 & \quad $z\gets z+1$ \\
 5 & \quad \textbf{if} $z\ge m$ \textbf{then} \\
 6 & \qquad \textbf{if} \textsc{PeerDemand}() \textbf{then} $B'\gets B+\Delta$ \\
 7 & \qquad $z\gets0$ \\
 8 & \textbf{else} $z\gets0$ \\
 9 & $B\gets\textsc{RoundDownPage}(\min(C,B'))$ \\
 10 & \textsc{PublishBudget}$(B)$; \textsc{EmitCapacityEvent}() \\
 \bottomrule
 \end{tabularx}
\end{minipage}\par\medskip

Starting at $B=z=0$, pressure (allocation stalls, reserve violations or
memory-related admission waits) reduces $B$ by $0<\gamma<1$. After $m$
stable windows below a lower wait threshold, deferred positive-value Peer
demand, assessed before quota filtering, permits an additive byte step
$\Delta$. Intermediate windows reset
$z$; missing telemetry disables growth. The ceiling clips every update;
emergencies may decrease within a window, but growth requires completed
stable windows. Parameters and window length are configurable.

Adapter changes neither fixed $Z$ nor per-request $L/K$. When occupancy
plus reservations exceeds the new quota, Master pauses admission and
selects safe relief. Physical leases must drain before freed bytes can
serve new work; insufficient reclaimable space causes backpressure.
The quota is authorization, not a claim that reclamation has completed.

\subsection{Connector: Committed Copies and Layered Recovery}
\label{sec:connector}

Connector uses bounded queues and service shares for foreground reuse,
recovery and checkpoint movement; real misses precede speculative reads
within the recovery share. Three invariants govern transfers:
\emph{complete compatible state before publication},
\emph{publish the replacement/fallback before retirement}, and
\emph{drain readers and copies before physical release}.
A DRAM-to-storage move therefore persists and commits the full manifest
before redirecting readers and releasing the source. Last-copy deletion
first publishes an older complete checkpoint or prompt-and-log fallback,
and withdraws unsupported Peer claims. Partial or stale callbacks cannot
advance $M$; failed copies preserve the prior authoritative path.

Recovery leases the last complete manifest on the bound target. Peer hits
bypass lower-tier reads/H2D; misses come from DRAM or storage then H2D.
Layer-ordered prefetch overlaps only remaining replay computation;
actual misses wait for all native KV, selector and TP dependencies.
Replay processes acknowledged tokens beyond $M$, suppresses duplicates
and resumes decode. Adopted Peer bytes move to execution accounting
without being freed or double-counted. These commit/release invariants
make prediction errors a performance issue rather than permission to
use missing state; progress still requires a compatible target and
available execution/I/O capacity.
